\section{Laplacetransformation}\label{MA2-laplacetransformation}

\begin{table}[H]
  \centering
  \begin{tabularx}{\textwidth}{@{}>{\raggedright\arraybackslash}X r@{}}
    \toprule
    \multicolumn{2}{c}{\textbf{�bersicht �ber Lernziele und Voraussetzungen}} \\
    \midrule
    \textbf{Lernziele} & \\[.5\normalbaselineskip]
    \tabitem die Laplace-Transformierte geeigneter Funktionen berechnen & \\
    \tabitem die Wirkung der Laplace-Transformation auf Ableitungen als Regel anwenden & \\[.5\normalbaselineskip]
    \midrule
    \textbf{Voraussetzungen} &  \\[.5\normalbaselineskip]
    \tabitem uneigentliche Integrale berechnen k�nnen & \abschnittsverweis{MA2-integration}\\
    \tabitem Ableitungen sicher berechnen k�nnen & \abschnittsverweis{MA1-ableitungsregeln} \\[.5\normalbaselineskip]
    \bottomrule
  \end{tabularx}
\end{table}

\noindent \ifintegriereaufgaben Die �bungsaufgaben zum Thema finden Sie \hyperref[MA2-laplacetransformation-aufgaben]{hier}.
\vspace{0.3cm}
\hrule
\fi

\subsection{Definition}

In vielen Anwendungen � insbesondere in der Regelungstechnik �
interessiert man sich f�r Vorg�nge, die zu einem bestimmten Zeitpunkt beginnen.
Ein System befindet sich im Ruhezustand und wird bei $t=0$
durch eine St�rung oder ein Eingangssignal angeregt.

Mathematisch f�hrt dies auf Anfangswertprobleme der Form
\[
	m \ddot x + \rho \dot x + D x = u(t),
	\qquad x(0)=x_0, \quad \dot x(0)=v_0.
\]

Die Schwierigkeit besteht darin, dass Ableitungen auftreten.
Ableitungen machen Gleichungen kompliziert �
Multiplikationen dagegen sind einfach.

\medskip

Die zentrale Idee der Laplace-Transformation lautet daher:

\begin{center}
	\emph{Man sucht eine Darstellung, in der Ableitungen zu Multiplikationen werden.}
\end{center}

Genau das leistet die Laplace-Transformation.
Sie ordnet einem Signal $y(t)$ eine Funktion $Y(s)$ zu,
in der Differentiation im Zeitbereich
zu einer algebraischen Operation im Bildraum wird.

\medskip

Dabei betrachtet man nur Zeiten $t \ge 0$.
Vor dem Zeitpunkt $t=0$ sei das System im Ruhezustand,
man setzt also meist $y(t)=0$ f�r $t<0$.

Die Laplace-Transformierte eines Signals
$y:[0,\infty)\to\mathbb C$ wird definiert als:
\begin{defbox}{Laplace-Transformation}{ }
	\[
		\mathcal L(y): D_Y \subset \mathbb C \to \mathbb C,
		\qquad
		s \mapsto Y(s),
		\qquad
		Y(s) := \int_0^\infty y(t)\, e^{-st}\,\mathrm d t.
	\]
\end{defbox}

\noindent Bemerkungen:
\begin{itemize}
	\item Das Integral muss nicht f�r alle $s\in\mathbb C$ existieren.
	Der Definitionsbereich $D_Y$ ist im Allgemeinen
	eine echte Teilmenge der komplexen Zahlen.

	\item Die Exponentialfunktion $e^{-st}$ wirkt f�r $\mathrm{Re}(s)>0$
	d�mpfend.
	W�chst $y(t)$ nicht schneller als eine Exponentialfunktion,
	so existiert das entstehende uneigentliche Integral.
\end{itemize}

\medskip

Man sagt: Das Signal $y(t)$ liegt im \emph{Zeitbereich}.
Seine Laplace-Transformierte $Y(s)$ liegt im \emph{Bildraum}.

Die Transformation ver�ndert nicht das physikalische System,
sondern nur seine Darstellung.
Der Perspektivwechsel macht jedoch eine entscheidende Struktur sichtbar:
Differentialgleichungen werden im Bildraum zu algebraischen Gleichungen.

\subsection*{Eigenschaften}

Was passiert, wenn man die Laplace-Transformation auf die Ableitung $\dot y$ eines Signals los l�sst?
\[
	\int\limits_0^\infty \dot y(t) \, e^{-st}\, \mathrm{d}t = \Big[ y(t) e^{-st} \Big]_0^\infty + s \int\limits_0^\infty y(t) \, e^{-st}\, \mathrm{d}t = -y(0) +s Y(s) \,.
\]
Die Ableitung wird also quasi ein Produkt - quasi, weil der Randterm noch als additive Konstante hinzukommt!\footnote{Die Fourier-Transformation der Ableitung liefert �brigens exakt ein Produkt.}\marginpar{\small Aufgabe}

In �hnlicher Weise kann man die Frage stellen, ob sich die Ableitung der Laplace-Transformierten in ein Produkt im Zeitbereich darstellt. Die Antwort ist ja. Die Formeln f�r h�here Ableitungen sind in der folgenden Box zusammengefasst:
\begin{theorembox}{Laplace-Transformierte und die Ableitung}{ }
	Die Laplace-Transformierte von \(y^{(n)}\) lautet:
	\[
		\mathcal L \left( y^{(n)}\right) (s) = s^n Y(s) - \sum\limits_{k=0}^{n-1} y^{(k)}(0) s^{n-k-1}
	\]
	Die $n$-te Ableitung der Laplace-Transformierten von \(y\) lautet:
	\[
		Y^{(n)}(s) = \mathcal L \left( (-t)^n y\right)(s)\,.
	\]
\end{theorembox}

\begin{exabox}{Laplace-Transformation von $\exp$, $\sin$ und $\cos$}{ }
	\begin{itemize}
	\item F�r $y:t\mapsto e^{\lambda t}$ mit $\lambda \in \mathbb R$ gilt
	\[
		\int\limits_0^\infty e^{\lambda t}\cdot e^{-st} \, \mathrm{d} t =
		\int\limits_0^\infty e^{ (\lambda-s) t} \, \mathrm{d} t =
	\Big[ \frac{e^{(\lambda -s ) t }}{\lambda -s}\Big]_0^\infty = \frac{1}{s-\lambda},\quad \textnormal{wenn } \mathrm{Re}(s)>\lambda\,.
	\]
	\item F�r $y:t\mapsto e^{i\omega t}$ mit $\omega \in \mathbb R$ gilt
	\[
		\int\limits_0^\infty e^{i\omega t}\cdot e^{-st} \, \mathrm{d} t =
		\int\limits_0^\infty e^{ (i\omega-s) t} \, \mathrm{d} t =
	\Big[ \frac{e^{(i\omega -s ) t }}{i\omega -s}\Big]_0^\infty = \frac{1}{s-i\omega},\quad \textnormal{wenn } \mathrm{Re}(s)>0\,.
	\]
\item F�r $y:t\mapsto \cos(\omega t)$ gilt
	\[
		\int\limits_0^\infty \cos(\omega t)\cdot e^{-st} \, \mathrm{d} t =
		  \frac{s}{s^2 + \omega^2},\quad \textnormal{wenn } \mathrm{Re}(s)>0\,.
	\]
\item F�r $y:t\mapsto \sin(\omega t)$ gilt
	\[
		\int\limits_0^\infty \sin(\omega t)\cdot e^{-st} \, \mathrm{d} t =
		  \frac{\omega}{s^2 + \omega^2},\quad \textnormal{wenn } \mathrm{Re}(s)>0\,.
	\]
	\end{itemize}
\end{exabox}

\begin{exabox}{Laplace-Transformation von Potenzfunktionen}{ }
Das Integral
\[
\int\limits_0^\infty t^n \cdot e^{-st}\,\mathrm{d}t  = \textnormal{?}
\]

Um das Integral zu l�sen ist es geschickt, den Integranden so umzuschreiben, dass eine Ableitung nach \(s\) auftaucht:
\begin{eqnarray*}
	\int\limits_0^\infty t^n \cdot e^{-st}\,\mathrm{d}t  &=& \int\limits_0^\infty (-1)^n\cdot \frac{\mathrm{d}^n}{\mathrm{d}s^n} \left( e^{-st} \right) \,\mathrm{d}t = (-1)^n\cdot \frac{\mathrm{d}^n}{\mathrm{d}s^n} \int\limits_0^\infty  e^{-st} \,\mathrm{d}t = (-1)^n\cdot \frac{\mathrm{d}^n}{\mathrm{d}s^n} \Big[ \frac{e^{-st}}{-s} \Big]_0^{\infty}\\ &=& (-1)^n \cdot \frac{\mathrm{d}^n }{\mathrm{d}s^n}\frac{1}{s}, \quad \mathrm{Re}(s) > 0\,.
\end{eqnarray*}

%Die Potenzfunktionen  mit  Exponent  $0,1,2,3,\dots$ haben  also  diese  Laplace-Transformierten:
%\begin{center}
%	\begin{tabular}{c |c | c}
%		$n$ &  $y(t)$ & $Y(s)$\\
%		\hline
%		$0$ & $1$ & $s^{-1}$ \\
%		$1$ & $t$ & $s^{-2}$ \\
%		$2$ & $t^2$ & $2 s^{-3}$ \\
%		$3$ & $t^3$ & $6 s^{-4}$ \\
%		$\vdots$ & $\vdots$ & $\vdots$ \\
%		$n$ & $t^n$ & $n! s^{-n-1}$ \\
%	\end{tabular}
%\end{center}
\end{exabox}

\begin{theorembox}{Laplace-Transformierten der Potenzfunktionen}{}
	F�r die Potenzfunktion $y: t\mapsto t^n$ mit $n\in \mathbb N^+$ existiert die Laplace-Transformierte f�r alle \(s\in \mathbb C\) mit \(\mathrm{Re}(s) >0\) und sie lautet:
\[ Y: s\mapsto n!\, \frac{1}{s^{n+1}}\,.\]
\end{theorembox}

\begin{exabox}{Zeitliche Verz�gerung}{ }
Wenn  das  Signal $y$ um  die  Zeit $a$ verz�gert  wird,  ergibt  sich  als  Laplace-Transformierte:
\[
	\int\limits_{\textcolor{blue}{a}}^\infty y(t\textcolor{blue}{-a}) e^{\displaystyle{-st}}\,\mathrm{d}t \stackrel{(*)}{=}
	\int\limits_0^\infty y(u) e^{\displaystyle{-s(u+\textcolor{blue}{a})}}\,\mathrm{d}u = e^{\displaystyle{-s\textcolor{blue}{a}}}\left( \mathcal L y\right) (s)\,.
\]
$(*)$ mit der Substitution $u=t-\textcolor{blue}{a}$, $\mathrm{d}u = \mathrm{d}t$.
\end{exabox}

\begin{exabox}{Zeitliche Beschleunigung}{ }
Wenn das Signal $y$ um den Faktor $b > 1$ beschleunigt (oder bei $b<1$ verlangsamt) wird, ergibt sich als Laplace-Transformierte:
\[
	\int\limits_{0}^\infty y(\textcolor{blue}{b}\cdot t) e^{\displaystyle{-st}}\,\mathrm{d}t \stackrel{(*)}{=} \int\limits_0^\infty y(u) e^{\displaystyle{ -\dfrac{s}{\textcolor{blue}{b}} \cdot u}} \frac{1}{\textcolor{blue}{b}}\,\mathrm{d}u = \frac{1}{\textcolor{blue}{b}}\left( \mathcal L y\right) \left( \dfrac{ s}{\textcolor{blue}{b}}\right)\,.
\]
$(*)$ mit der Substitution $u=\textcolor{blue}{b} \cdot t$, $\mathrm{d}u = \textcolor{blue}{b}\,\mathrm{d}t$.
\end{exabox}

\begin{theorembox}{Eigenschaften der Laplace-Transformation}{ }
	Die Laplace-Transformation $\mathcal L: f\mapsto \mathcal L(f)$ ist eine lineare Abblidung:
	\begin{itemize}
		\item $\mathcal L(f+g) = \mathcal L(f) + \mathcal L(g) $
		\item $\mathcal L(\lambda f) = \lambda \mathcal L(f)$
	\end{itemize}
\end{theorembox}

F�r eine kausale Verz�gerung verwenden wir die Heaviside-Funktion
\(H(t)=0\) f�r \(t<0\) und \(H(t)=1\) f�r \(t\geq0\).
Das um \(a>0\) verz�gerte Signal lautet damit \(H(t-a)f(t-a)\).

\begin{exabox}{Rechenregeln f�r die Laplace-Transformation}{ }
	Die Laplacetransformierten der Standardsignale \(f_1: t\mapsto \cos(t)\) und \(f_2: t\mapsto t\) sind gegeben:
	\[ \cos(t) \,\laplace \, \dfrac{s}{s^2+1}, \quad t \,\laplace \, \dfrac{1}{s^2}\]
	Mit Hilfe der Rechenregeln folgt:
	\[
		\begin{array}{l rclcl l}
			f:t\mapsto \cos(t) + t\; \quad & \mathcal L(f)(s) &=& \mathcal L(f_1)(s) + \mathcal L (f_2) (s) &=& \dfrac{s}{s^2+1} +\dfrac{1}{s^2}  & \textnormal{ Linearitaet} \\[2ex]
			f:t\mapsto H(t-3)\cos(t-3)\; \quad & \mathcal L(f)(s) &=& e^{-3\, s} \cdot \mathcal L(f_1)(s) &=& e^{-3\, s}\dfrac{s}{s^2+1} & \textnormal{ Verz�gerung} \\[2ex]
			f:t\mapsto \cos(4t)\; \quad & \mathcal L(f)(s) &=& \dfrac14 \,\mathcal L (f_1)\left( \dfrac{s}{4} \right) &=& \dfrac14 \dfrac{\dfrac{s}{4}}{\left( \dfrac{s}{4}\right)^2+1} = \dfrac{s}{s^2+16}  & \textnormal{ Beschleunigung} \\[2ex]
			f:t\mapsto \cos(4t) - t\; \quad & \mathcal L(f)(s) &=&  \dfrac14 \, \mathcal L(f_1)\left( \dfrac{s}{4}\right)- \mathcal L (f_2)(s)  &=& \dfrac{s}{s^2+16}-\dfrac{1}{s^2} & \textnormal{ gemischt}
		\end{array}
	\]
\end{exabox}


Es gibt eine f�r die Praxis wichtige Eigenschaft der Laplacetransformation. Um diese Eigenschaft aufzuschreiben, definieren wir die sogenannte Faltung:

\begin{defbox}{ Faltung }{ }
Seien \(f,g: \mathbb R \to \mathbb K\). Die Faltung von \(f\) und \(g\) ist die folgende Funktion, falls sie existiert: \[ \left( f\star g\right)(t) = \int\limits_{-\infty}^{\infty} f(t-\tau)\cdot g(\tau)\,\mathrm{d}\tau\,. \] Es gilt

\begin{itemize}
	\item[(a)] \( f\star(g\star h) = (f\star g) \star h\)
    \item[(b)]  \( f\star (g+h) = f\star g + f\star h\)
    \item[(c)] \(f\star g = g\star f\)
    \item[(d)] \( (\lambda f)\star g = \lambda \cdot (f\star g), \quad \lambda \in \mathbb K\)
\end{itemize}

Gilt \( f(x) = 0\) und \( g(x)=0\) f�r \(x<0\) so folgt  \[ (f\star g)(x) = \int\limits_0^x f(x-\xi) g(\xi)\,\mathrm{d}\xi \,.\]
\end{defbox}

Die Faltung ist ein Integraloperator, der assoziativ, distributiv und kommutativ ist und im Kontext der Laplace-Transformation gilt folgendes:
2. Anwendung im Kontext der Laplace-Transformation

\begin{theorembox}{ Laplace-Transformation einer Faltung}{ }
Seien \(f,g:[0,\infty) \to \mathbb R\) st�ckweise stetig mit Laplace-Transformationen \(F,G\). Dann gilt f�r die Faltung: \[ \left( f\star g\right)(t) = \int\limits_0^t f(t-\tau)\cdot g(\tau)\,\mathrm{d}\tau \] die Produktformel \[ \mathcal L\left( f\star g\right)(s) = \mathcal L(f)(s) \cdot \mathcal L(g)(s)\,. \]
\end{theorembox}

%\subsection{Zusammenfassung �ber Transformationen}
%
%Damit kennen wir nun die vier wichtigsten Funktionaltransformationen:\marginpar{\small Aufgabe}
%\begin{mybox}[Transformationen von Funktionen]
%	\begin{center}
%		\begin{tabular}{l ||l|l}
%			Transforamtion & Fourier-Reihe und DFT (FFT) & $z$-Transformation  \\
%			Anwendung & Signalverarbeitung  & Signalverarbeitung \\
%			\hline
%			Transformation & (kontinuierliche) Fourier-Transformation & Laplace-Transformation \\
%			Anwendung & Optik, Akustik, Elektrodynamik  & Regelungsprozesse
%		\end{tabular}
%	\end{center}
%\end{mybox}

